\title{Parameter-Efficient Orthogonal Finetuning via Factorization}

\begin{abstract}
The proliferation of large foundation models has made them increasingly prevalent; however, training these models from the ground up incurs substantial computational cost. Consequently, devising methods that facilitate the efficient adaptation of these robust models to specific downstream tasks is of growing significance. This work investigates a systematic approach to finetuning, namely Orthogonal Finetuning (OFT), to enhance task-specific model adaptation. While OFT provides strong generalization, it typically requires a sizable number of trainable parameters, owing to the inherently high dimensionality associated with orthogonal matrices. To mitigate this, we analyze OFT through the lens of information transmission, extracting critical criteria that promote greater parameter efficiency. Drawing inspiration from the Cooley-Tukey fast Fourier transform’s capability for streamlined information transfer, we introduce a more compact orthogonal parameterization utilizing butterfly structures. By integrating this structure within OFT, we present a new, parameter-efficient finetuning technique termed Orthogonal Butterfly~(BOFT). BOFT generalizes OFT, encompassing it as a particular case and establishing a broad orthogonal finetuning paradigm. We further undertake comprehensive experiments, examining the adaptation of large vision transformers, substantial language models, and text-to-image diffusion models for a spectrum of vision and language downstream applications.
\end{abstract}

\section{Introduction}

The advent of models such as ChatGPT~\cite{brown2020language,chen2021evaluating} and Stable Diffusion~\cite{rombach2022high} exemplifies the exceptional generalization capacities achieved by large-scale foundation models. Accompanying these advances is a significant growth in parameter count—

(\emph{e.g.}, GPT-3 boasts roughly 175B parameters)—which escalates the difficulty and resource requirements of training models from the outset. As such, continued progress increasingly hinges on the capacity to repurpose pre-existing foundation models for novel tasks, rather than embarking on exhaustive training cycles. Effectively, it is essential to adapt existing foundation models to new domains with maximal efficiency. Approaches for task adaptation predominantly fall into three categories: \emph{model finetuning}~\cite{hulora2022,zaken2022bitfit,guo2020parameter,raffel2020exploring,qiu2023controlling,zhang2022adaptive,chavan2023one}, where only certain parameters are updated while retaining the original model structure; \emph{adapter tuning}~\cite{rebuffi2017learning,houlsby2019parameter,pfeiffer2020adapterfusion,lin2020exploring,he2021towards}, which supplements the model with new, trainable modules; and \emph{prompt tuning}~\cite{li2021prefix,lester2021power}, wherein trainable prefix tokens are appended to the input sequence. Among these, model finetuning is particularly notable for its conceptual simplicity, effectiveness, and lack of added inference overhead.

Model finetuning fundamentally aims to maintain a certain degree of similarity between the pretrained and finetuned models, as quantified by specific metrics, thereby retaining knowledge acquired during pretraining. Commonly, this is achieved by employing a relatively low learning rate, an intuitive strategy that limits divergence between the pretrained and the updated models. Recognizing the inherent structure of weight matrices, more systematic strategies attempt to retain the internal relationships within these matrices, particularly by preserving the pairwise angular relationships among neurons. This motivation underpins the design of Orthogonal Finetuning~(OFT)~\cite{qiu2023controlling}, a recent framework where every neuron within a layer is subject to a shared orthogonal transformation, thereby guaranteeing the invariance of pairwise angles during finetuning. While OFT has exhibited favorable generalization and convergence properties in the context of text-to-image diffusion model adaptation~\cite{qiu2023controlling}, it faces a significant drawback: the number of parameters grows rapidly with the dimensionality of the orthogonal matrices involved. To mitigate this, OFT proposes using a block-diagonal orthogonal structure to cut down on parameter count. Nevertheless, this comes with certain trade-offs—the imposed block-diagonal sparsity is rigid and results in independent transformations across blocks. Such a sparsity scheme can lead to adverse inductive biases (for instance, limiting the expressive power of the transformation or failing to capture standard linear mappings), and, crucially, the challenge of selecting an optimal sparsity structure remains unresolved.

The central challenge lies in constructing a dense orthogonal matrix in a manner that conserves parameters. At first, this appears impractical, as a fully dense orthogonal matrix in $d$ dimensions inherently demands $\mathcal{O}(d^2)$ parameters. To overcome this, we introduce a method that synthesizes a dense orthogonal matrix by composing multiple factorized sparse matrices. This technique is rooted in the realization that one can decrease the number of trainable parameters by exchanging storage requirements for computational overhead: since the orthogonal matrix is encoded as a sequence of sparse matrix multiplications, these products must be recomputed at every training step. To clarify the role of this matrix factorization, we recast the generation of a dense orthogonal matrix as an exercise in information transmission. In this perspective, constructing a dense orthogonal matrix through the multiplication of sparse factors corresponds to conveying information across a grid-based graph. This conceptualization paves the way for designing a multitude of sparse matrix factorizations that not only constrain the parameter count but also retain sufficient capacity to model dense transformations.

To realize parameter savings within this information transmission paradigm, we take cues from the butterfly architecture employed in the Cooley-Tukey fast Fourier transform algorithm~\cite{cooley1965algorithm}, which enables highly efficient communication between graph nodes~\cite{klasing1994broadcasting}. The butterfly connectivity pattern implicit in the Cooley-Tukey method suggests a specific approach for sparse matrix factorization, referred to as butterfly factorization. With butterfly factorization, it becomes feasible to represent a dense $d$-dimensional matrix as a product of $\mathcal{O}(\log d)$ sparse matrices, with each factor containing just $\mathcal{O}(d)$ nonzero elements. By maintaining orthogonality in each sparse factor, this leads to a parameter-efficient orthogonal representation that requires only $\mathcal{O}(d\log d)$ parameters—a dramatic reduction compared to the conventional approach. Capitalizing on this compact orthogonal parameterization, we introduce a novel and efficient finetuning strategy, termed Orthogonal Butterfly~(BOFT). BOFT generalizes OFT as a particular case, offering a flexible framework for orthogonal finetuning. Both block-diagonal and butterfly arrangements share the attribute of sparsity; in each, imposing sparsity on the orthogonal matrices diminishes the number of trainable parameters required. When comparing our strategy to the low-rank decompositions employed in LoRA~\cite{hulora2022}, it is notable that both low-rank and sparse matrices form prominent categories of structured matrices~\cite{candes2011robust} aimed at achieving parameter efficiency.

In contrast to OFT, which balances expressivity and regularity through a block-diagonal parameterization, BOFT leverages the butterfly structure to achieve a more continuous spectrum between fully expressive matrices from the orthogonal group and the identity matrix associated with regularity. This approach allows us to define a narrower, yet sufficiently expressive, subset of orthogonal transformations suitable for downstream applications. The butterfly structure, a foundational component in efficient linear operations like the discrete Fourier and cosine transforms, motivates our argument that adopting such parameter-efficient architectures provides an implicit inductive bias. This bias promotes generalization and mitigates overfitting risks. Notably, while the butterfly structure is capable of representing a variety of well-known linear transforms, the block-diagonal approach utilized in OFT cannot do so. Our primary contributions are summarized below:

\begin{itemize}[leftmargin=*,nosep]
    \item We frame the challenge of parameter efficiency in orthogonal finetuning through a novel information transmission perspective. Within this framework, the design of a dense, parameter-efficient orthogonal matrix is recast as an information flow problem on a grid-shaped graph.
    \item Drawing on the butterfly configuration characteristic of the Cooley-Tukey algorithm, we introduce Orthogonal Butterfly, an innovative and compact orthogonal finetuning strategy based on the information transmission viewpoint.
    \item We offer several theoretical justifications for how BOFT can drastically decrease the parameter count, yet preserve strong expressivity and training robustness. Additionally, the factorized matrix structure in BOFT naturally supports a unique form of weight interpolation (refer to Figure~\ref{fig:boft_interpolation}).
    \item We are the first to demonstrate the effectiveness of orthogonal finetuning across diverse applications beyond controllable text-to-image synthesis~\cite{qiu2023controlling}. We empirically evaluate BOFT on adaptation challenges in domains from computer vision to natural language processing, showing that it significantly outperforms existing state-of-the-art methods, thereby affirming its efficiency and broad applicability.
\end{itemize}

\section{Related Work}

\textbf{Parameter-efficient finetuning (PEFT)}. As foundation models (\emph{e.g.}, \cite{brown2020language,radford2021learning,rombach2022high,kirillov2023segment}) grow larger and more capable, the challenge of efficiently adapting these models to new tasks has received growing attention~\cite{treviso2023efficient,ding2023parameter,lialin2023scaling}. Numerous PEFT strategies have emerged~\cite{houlsby2019parameter,aghajanyan2020intrinsic,hulora2022,edalati2022krona,wang2022adamix,gheini2021cross,zaken2022bitfit,guo2020parameter,sung2021training,ansell2022composable,lester2021power,li2021prefix,vu2022spot,he2021towards,mao2021unipelt,karimi2021compacter,liu2022few,sung2022lst,chen2022parameter,jia2022visual,chen2022adaptformer,zhang2022neural,jie2023fact,lian2022scaling,luo2023towards}, with reparameterization-centric techniques~\cite{aghajanyan2020intrinsic,hulora2022,edalati2022krona} being particularly pertinent to our study. Among these, LoRA~\cite{hulora2022} modifies pretrained weights by introducing a sum of two low-rank matrices, demonstrating strong results for language tasks. Since LoRA fixes the rank of each added matrix uniformly across layers, subsequent efforts~\cite{zhang2022adaptive,valipour2022dylora,zhang2023increlora} propose adapting matrix rank per layer to better allocate parameters under a given budget. Owing to its conceptual simplicity, this form of low-rank reparameterization is now widely adopted~\cite{dettmers2023qlora,zi2023delta,chavan2023one}. Drawing from insights on hyperspherical energy and generalization~\cite{liu2018learning,liu2021orthogonal}, orthogonal finetuning (OFT)~\cite{qiu2023controlling} emerges as an alternative, targeting diffusion models for text-to-image generation. OFT optimizes an orthogonal transformation for each layer's neurons, yielding improved generalization and more stable training dynamics than LoRA, albeit at the cost of increasing the number of trainable parameters. Reducing OFT’s parameter count, therefore, remains an open and valuable research aim. In addition, whether the benefits of OFT generalize beyond specific domains (such as text-to-image diffusion) remains an open question. BOFT seeks to close this gap by employing butterfly factorization to boost the parameter-efficiency of OFT. With this, the promise of orthogonal finetuning is extended to broader adaptation scenarios.

\textbf{Butterfly structures}. The radix-2 Cooley-Tukey procedure recursively decomposes the $N$-point discrete Fourier transform into two $\frac{N}{2}$-point discrete Fourier transforms, yielding what is termed a butterfly structure—a product of several sparse matrices, collectively referred to as a butterfly matrix. This type of matrix has proven useful in diverse settings, including orthogonal matrix parameterization to bypass pivoting in Gaussian elimination~\cite{parker1995random}, stabilizing recurrent neural network training~\cite{jing2017tunable}, and kernel approximation~\cite{munkhoeva2018quadrature}. Additionally, works such as \cite{dao2019learning,dao2019kaleidoscope} develop efficient linear transforms via butterfly matrix representations, while \cite{chen2022pixelated,dao2022monarch} show the utility of butterfly matrices for scaling up neural network training. Beyond these, butterfly structures facilitate rapid matrix-vector computations~\cite{michielssen1996multilevel,o2010algorithm}, support efficient data-sparse matrix approximations~\cite{li2015butterfly}, and assist with information transmission in networks~\cite{klasing1994broadcasting,li2003linear,rajasingh2016transmission}. Our approach diverges from earlier work by leveraging butterfly structures to render OFT more parameter-efficient.

\section{An Information Transmission View on Orthogonal Finetuning}

We begin by outlining foundational aspects of OFT. In order to adapt the pretrained weight matrix, OFT introduces a reparameterization where the new weight matrix is expressed as the product of a trainable orthogonal matrix and the fixed pretrained matrix. This differs from LoRA, which employs an additive update using a low-rank matrix; OFT, by contrast, utilizes a multiplicative approach via an orthogonal matrix. For parameter efficiency, LoRA adopts a low-rank design, while the original OFT instead opts for a block-diagonal orthogonal structure~\cite{qiu2023controlling}. Notably, parameter efficiency in OFT improves as the block size along the diagonal decreases. Figure~\ref{fig:comp_lora} provides a visual summary of these distinctions. 

The primary rationale for employing an orthogonal transformation when finetuning the weights is to maintain the pairwise angles between neuron activations~\cite{liu2018learning,liu2021orthogonal,qiu2023controlling}, thereby retaining a substantial amount of the pretrained model’s semantic information. Specifically, OFT learns an orthogonal matrix $\bm{R}\in\mathbb{R}^{d\times d}$ for a pretrained linear mapping $\bm{W}^0\in\mathbb{R}^{d\times n}$, modifying the layer's forward computation from $\bm{z}=(\bm{W}^0)^\top\bm{x}$ to $\bm{z}=(\bm{R}\bm{W}^0)^\top\bm{x}$, where $\bm{x}\in\mathbb{R}^{d}$ denotes the input vector and $\bm{z}\in\mathbb{R}^{n}$ the output. To realize a form of zero initialization, OFT sets $\bm{R}$ as the identity matrix at the outset. Preserving the orthogonality of $\bm{R}$ during finetuning relies on the Cayley parameterization~\cite{qiu2023controlling,liu2021orthogonal}: $\bm{R}=(\bm{I}+\bm{Q})(\bm{I}-\bm{Q})^{-1}$, where $\bm{Q}$ is a skew-symmetric matrix (i.e., $\bm{Q}=-\bm{Q}^\top$).

To further enhance parameter efficiency, OFT models $\bm{R}$ with a block-diagonal structure, where $\bm{R}=\text{diag}(\bm{R}_1,\bm{R}_2,\cdots,\bm{R}_r)$ and each $\bm{R}_i\in\mathcal{R}^{b\times b}$ is a small orthogonal block, with the total dimensionality satisfying $br=d$. While this configuration reduces parameter count, it imposes the constraint that each neuron's dimensions (corresponding to each column of $\bm{W}^0$) are split into $r$ groups, each manipulated independently by its own orthogonal matrix. Although this block-diagonal method demonstrates practical merit, the arbitrary division of neuron dimensions solely by their indices lacks strong justification, motivating interest in dense orthogonal matrices. This naturally leads to the question: \emph{Is it feasible to construct a dense orthogonal matrix that maintains parameter efficiency?}

To tackle this issue, we introduce a method where a dense orthogonal matrix is decomposed as a product of several sparse orthogonal matrices. To facilitate a systematic and comprehensible analysis of the sparsity structures within orthogonal matrix decomposition, we reinterpret the generation of a dense orthogonal matrix in OFT through the lens of information transmission. Specifically, constructing a dense matrix $\bm{R}\in\mathbb{R}^{d\times d}$ via the multiplication of $m$ square matrices, $\bm{R} = \bm{B}_m\bm{B}_{m-1}\cdots\bm{B}_1$, is analogous to transmitting information across a lattice of $d\times (m+1)$ nodes, as shown in Figure~\ref{fig:info_trans}. The rationale for this perspective stems from the fact that a dense $d$-by-$d$ matrix naturally models dense connections from one set of $d$ nodes to another. For instance, a zero entry $R_{ij}$ in $\bm{R}$ signals that communication from node $j$ to node $i$ is blocked, whereas a non-zero value allows such a pathway. Thus, expressing $\bm{R}$ as $\bm{B}_m\bm{B}_{m-1}\cdots\bm{B}_1$ can be viewed as successive exchanges of information, dictated by the connectivity defined by each $\bm{B}_i$. The transmission sequence starts with $\bm{B}_1$ and proceeds through to $\bm{B}_n$. As a concrete case in Figure~\ref{fig:info_trans}, consider the decomposition $\bm{R} = \bm{B}_5\bm{B}_4\bm{B}_3\bm{B}_2\bm{B}_1$; both the corresponding sparsity and the induced graphical structure are visualized. This graphical depiction arises by unraveling the chain of matrix multiplications, where each $\bm{B}_i$ defines directed links from level $i$ to level $(i+1)$ nodes. More precisely, the $(j_1,j_2)$-th entry of $\bm{B}_i$ specifies whether a directed connection exists from the $j_2$-th node in level $i$ to the $j_1$-th node in level $(i+1)$ (with a zero entry denoting the absence of an edge). In order for $\bm{B}_5\bm{B}_4\bm{B}_3\bm{B}_2\bm{B}_1$ to produce a dense matrix, it is required that each initial node in the first layer can propagate information to all nodes in the sixth layer. Looking at $\bm{R} = \bm{B}_4\bm{B}_3\bm{B}_2\bm{B}_1$, representing communication from sources in the first layer to receivers in the fifth, it is apparent that node $1$ is unable to transmit to node $3$; consequently, the $(3,1)$ position in the sparsity pattern of $\bm{B}_4\bm{B}_3\bm{B}_2\bm{B}_1$ is zero. Similar relationships between the induced graph and the nonzero patterns are evident for $\bm{R} = \bm{B}_3\bm{B}_2\bm{B}_1$ and $\bm{R} = \bm{B}_2\bm{B}_1$. More generally, for any dense matrix $\bm{R}\in\mathbb{R}^{d\times d}$, the product $\bm{B}_m,\cdots,\bm{B}_1$ must map to a configuration of directed connections on a $d\times (m+1)$ grid, with each edge linking nodes from neighboring columns (i.e., adjacent levels), ensuring that every node at the initial level has a transmission path to all nodes at the final level $(m+1)$.

The information transmission perspective offers valuable insights into the sparsity configurations of factorization matrices utilized in OFT. In Figure~\ref{fig:blk_diag}, the block-diagonal layout of $\bm{R}$ within the original OFT is depicted. Although this design choice successfully limits the quantity of trainable parameters, it falls short of generating a fully dense matrix $\bm{R}$. Our objective is to synthesize a dense orthogonal matrix from $m$ sparse orthogonal factors, while minimizing the count of trainable parameters that effectively contribute. Interpreted through the information transmission lens, our primary criteria are: \emph{(i)} \textbf{dense connectivity}, ensuring that each node at the initial layer can reach every node at the final layer via at least one path; and \emph{(ii)} \textbf{minimum free edges}, seeking the lowest possible edge count compatible with the orthogonality requirement. The orthogonality condition imposes a nuanced limitation on the potential connections (edges) between adjacent layers. For instance, each matrix $\bm{B}_i$ must be of full rank to uphold orthogonality, which necessitates $d$ connecting edges that establish a one-to-one mapping between nodes in the $i$-th and $(i-1)$-th levels—yielding at least $d$ such connections (\emph{e.g.}, four in the case illustrated by Figure~\ref{fig:info_trans}). These $d$ essential connections are mandated by orthogonality and are not considered part of the edge budget, as they involve fixed, non-trainable parameters (\emph{e.g.}, in a $d \times d$ orthogonal matrix, $d$ non-zero elements can be only $\pm 1$). The constraint of orthogonality reduces the number of independent parameters relative to a full matrix, introducing interdependencies among the remaining edges. For example, a $2 \times 2$ orthogonal matrix with a zero at the $(1,1)$ entry necessarily has a corresponding zero at the $(2,2)$ entry; thus, eliminating one edge can dictate the removal of another. Therefore, the feasible patterns of connectivity may be altered by orthogonality, which can sometimes increase or decrease the number of permissible edges. Examining the pattern of nonzero entries in the resulting orthogonal matrix, the information transmission framework becomes especially apt for OFT, since our primary concern lies in the locations of trainable, nonzero elements in $\bm{R}$, and not in their exact values.

A straightforward approach to connect two layers densely requires $\mathcal{O}(d^2)$ connections (i.e., one full orthogonal matrix), amounting to $d^2-d$ connections—for instance, when $d=4$, this corresponds to 12 edges. As depicted in Figure~\ref{fig:info_trans}, an example of valid matrix factorization employs only $10$ edges in total, which is actually fewer than a fully dense orthogonal connection. By conceptualizing this framework graphically, it becomes possible to analyze the parameter efficiency of OFT and to readily devise plausible factorizations. This study draws on the structure of butterfly graphs~\cite{cooley1965algorithm}, a notable topology from the Cooley-Tukey algorithm, which achieves efficient full connectivity between $d$ sources and $d$ targets with just $\mathcal{O}(d\log d)$ edges. In fact, whereas the network structure illustrated in Figure~\ref{fig:info_trans} requires $10$ edges for dense connectivity, the butterfly graph accomplishes the same using only $8$ edges. The next section details how the butterfly configuration promotes greater parameter efficiency.

\section{Orthogonal Parameterization by Butterfly Factorization}

Originally, the butterfly pattern was used within the Cooley-Tukey algorithm to expedite the fast Fourier transform. Within the context of the Fourier transform, any localized modification in the frequency domain can affect the entire spatial domain—a phenomenon that mirrors the nature of our information transmission scenario, where each node in an initial layer can relay data to every node in the output layer. Additionally, butterfly topologies have been widely adopted in computer network designs~\cite{solihin2015fundamentals,li2011linear} due to their efficiency in facilitating information flow. Suppose $k\geq 2$ is a power of $2$. We define the \emph{butterfly factor} $\bm{B}^F(k)$ as follows:

\begin{equation}\label{eq:but_factor}
\begin{aligned}
    \bm{B}^F(k)=\begin{bmatrix}
    \text{diag}(\bm{d}_1) & \text{diag}(\bm{d}_2) \\
    \text{diag}(\bm{d}_3) & \text{diag}(\bm{d}_4)
    \end{bmatrix}\in\mathbb{R}^{k\times k},
\end{aligned}
\end{equation}

where each $\bm{d}_i\in\mathbb{R}^{\frac{k}{2}}$ for $i=1,2,3,4$ represent arbitrary vectors. Given $d=2^N$, the $d$-dimensional \emph{butterfly matrix} $\bm{B}(d)\in\mathbb{R}^{d\times d}$ is then recursively constructed as:

\begin{equation}
\begin{aligned}
    \!\!\bm{B}(d)=\tilde{\bm{B}}(d,d)\!\cdot\!\begin{bmatrix}
    \bm{B}_1(\frac{d}{2})\!\!\!\!\! & \bm{0} \\
    \bm{0}\!\!\!\!\! &\bm{B}_2(\frac{d}{2})
    \end{bmatrix}= \tilde{\bm{B}}(d,d)\tilde{\bm{B}}(d,\frac{d}{2})\cdots\tilde{\bm{B}}(d,2),
\end{aligned}
\end{equation}

Here, $\bm{B}_1(\frac{d}{2})$ and $\bm{B}_2(\frac{d}{2})$ represent two butterfly matrices, each of dimension $\frac{d}{2}$. We introduce the notion of a \emph{butterfly component} as $\thickmuskip=2mu \medmuskip=2mu \tilde{\bm{B}}(d,k)=\text{diag}(\bm{B}^F_1(k),\cdots,\bm{B}^F_{d/k}(k))$, which constitutes a block-diagonal matrix of dimensions $d\times d$ whose block size is $k$. Here, every $\bm{B}^F_i(k)$, for all $i$, denotes a butterfly factor specified by Equation~\ref{eq:but_factor}. This construction enables us to employ the butterfly matrix as a parameterization mechanism for orthogonal matrices. To facilitate this, it suffices to ensure the orthogonality of every block matrix $\tilde{\bm{B}}(d,k)$ within the butterfly matrix $\bm{B}(d)$ for each $k$. Focusing initially on the block-diagonal case with block size $2$, namely $\tilde{\bm{B}}(d,2)$, orthogonality can be imposed through the Cayley transform or $2$-dimensional rotation for each $2\times2$ block, consistent with \cite{liu2021orthogonal,qiu2023controlling}. Each butterfly component’s nonzero structure is a permutation of the pattern found in $\tilde{\bm{B}}(d,2)$, which ensures that all butterfly components admit straightforward parameterization as orthogonal matrices. This results in a highly efficient orthogonal matrix parameterization leveraging numerous $2\times2$ orthogonal blocks. Expanding on this idea, following \cite{chen2022pixelated}, we define a block butterfly component $\tilde{\bm{B}}^b(d,k)$ where every entry $\bm{d}_i$, for each $i$, is now a $b\times b$ matrix. To maintain the orthogonality of $\tilde{\bm{B}}^b(d,2)$, each $2b\times2b$ block must be parameterized as an orthogonal matrix. For the other butterfly components $\tilde{\bm{B}}^b(d,k)$ where $k>2$, the block-wise permutation of $\tilde{\bm{B}}^b(d,2)$'s sparsity pattern permits them to likewise be structured as orthogonal matrices. Bringing these components together, the forward computation in BOFT is expressed as

\begin{equation}
    \begin{aligned}\nonumber
    \bm{z}=\big(\bm{R}(m,b)\cdot\bm{W}^0\big)^\top\bm{x},~~~~\text{s.t.}~\bigg{\{}\bm{R}(m,b)=\prod_{i=1}^m \tilde{\bm{B}}^b_{(d,i)}~~\&~~\underbrace{\big(\tilde{\bm{B}}^b_{(d,j)}\big)^\top\tilde{\bm{B}}^b_{(d,j)}=\tilde{\bm{B}}^b_{(d,j)}\big(\tilde{\bm{B}}^b_{(d,j)}\big)^\top\!=\bm{I}_d}_{\forall j\in [1, m]}\bigg{\}},
    \end{aligned}
\end{equation}

Here, for notational simplicity, we let $\tilde{\bm{B}}^b(d,2^{m - i+1})$ be represented as $\tilde{\bm{B}}^b_{(d,i)}$, and $\bm{I}_d$ denotes the identity matrix of dimension $d$. The orthogonal transformation matrix $\bm{R}(m,b)\in\mathbb{R}^{d\times d}$ is structured as a product of $m$ orthogonal butterfly components. We use the notation BOFT with $\bm{R}(m,\frac{b}{2})$ as BOFT($m$,$b$) for clarity, requiring $b\geq 2$. When setting $m=1$, BOFT($1$,$b$) collapses to the block-diagonal OFT~\cite{qiu2023controlling}, where the block size is $b$. Choosing BOFT($1$,$d$) returns the vanilla OFT~\cite{qiu2023controlling} that employs a full orthogonal matrix without restriction. Moreover, selecting BOFT($\log \frac{2d}{b}$,$b$) corresponds to using the block butterfly matrix $\bm{B}^{\frac{b}{2}}(d)$ for $\bm{R}$, which leads to a densely populated orthogonal matrix. Broadly, BOFT($m$,$b$) incorporates $\frac{1}{2}(b-1)dm$ trainable parameters when adapting a $d\times n$ linear layer. When specialized to the butterfly matrix with $m=\log d, b=2$, BOFT only uses $\mathcal{O}(d\log d)$ parameters. In comparison, original OFT involving an unstructured orthogonal matrix consumes $\mathcal{O}(d^2)$ parameters, while block-diagonal OFT, with block count $r$, uses $\mathcal{O}(bd)$. Hence, to realize a dense orthogonal structure, standard OFT necessitates setting $b=d$, but BOFT attains such structure for any choice of $b$.

\textbf{Identity Initialization for BOFT}. In transfer learning setups, finetuning customarily begins from the original pretrained weights to ensure minimal deviation from the pretrained solution. For instance, LoRA opts for zero initialization on the additional low-rank matrices. Analogously, BOFT sets all butterfly components initially to identity matrices (i.e., the Cayley transform’s skew-symmetric component starts at zero).

\textbf{Multiplicative Dropout for BOFT}. LoRA~\cite{hulora2022} enhances generalization by applying a Dropout module to its low-rank updates to mitigate overfitting. While the traditional Dropout~\cite{srivastava2014dropout} is readily applicable in LoRA, it is incompatible with BOFT’s multiplicative updating mechanism. To overcome this limitation, we introduce a multiplicative variant of Dropout tailored for BOFT: given that $\bm{R}(m,b)$ comprises $m$ butterfly factors, which can be arranged as block-diagonal matrices of size $2b\times 2b$, our approach randomly selects $p_1$ percent of butterfly factors and, within each, randomly replaces $p_2$ percent of their diagonal sub-blocks by identity matrices.

\section{Intriguing Insights and Discussions}

\textbf{Expressivity of BOFT}. Previous work has demonstrated that the butterfly matrix structure, combined with suitable permutations, is capable of exactly representing a variety of well-known fast linear transforms~\cite{dao2019learning,dao2019kaleidoscope}, such as the fast Fourier transform and the Hadamard transform. However, the extent to which the orthogonal butterfly matrix can approximate arbitrary orthogonal matrices has not been thoroughly examined. To empirically study this question, we perform an experiment in which a random dense orthogonal matrix~\cite{anderson1987generation} of dimension $512 \times 512$ is approximated using BOFT. Each data point in Figure~\ref{fig:para_eff} reflects an average over 10 different random seeds. On the vertical axis, we plot the approximation error; the horizontal axis shows the effective number of trainable parameters. Each curve (distinguished by color) corresponds to a BOFT configuration with a specific block size, and the leftmost marker on every curve gives the approximation error for BOFT($1$,$b$) (that is, the standard block-diagonal OFT using block size $b$). BOFT generally achieves higher parameter efficiency compared to OFT. For instance, BOFT($9$,$2$) demonstrates greater expressivity than BOFT($1$,$16$) while utilizing significantly fewer parameters. In general, using a smaller $b$ together with a larger $m$ improves parameter efficiency: BOFT($6$,$4$) achieves comparable approximation error to BOFT($2$,$16$) with far fewer parameters. This implies that the butterfly matrix captures a richer subset of the orthogonal group relative to the block-diagonal architecture, substantiating that BOFT provides greater expressivity than OFT when the block size is fixed.

\begin{theorem}[Expressivity of BOFT]\label{thm:exp_boft}
    BOFT surpasses OFT in expressivity for an identical block size. By stacking butterfly matrices, one can approximate any orthogonal matrix of size $d$ as follows: $\bm{B}_{d-1,1}(d)\bm{B}^\top_{d-1,2}(d)\cdots\bm{B}_{1,1}(d)\bm{B}^\top_{1,2}(d)$, with each $\bm{B}_{i,j}(d)$ denoting a butterfly matrix.
\end{theorem}

Theorem~\ref{thm:exp_boft} naturally leads to a generalized formulation of BOFT: an orthogonal matrix can be represented as $\thickmuskip=2mu \medmuskip=2mu \bm{R}^G(m_1,b_1,m_2,b_2,l)=\bm{R}_{l,1}(m_1,b_1)\bm{R}_{l,2}^T(m_2,b_2)\cdots\bm{R}_{1,1}(m_1,b_1)\bm{R}_{1,2}^T(m_2,b_2)$, where $\bm{R}_{i,j}^T(m,b)$ is the orthogonal matrix employed in the BOFT construction. In the special case where $m_1=m_2=\log d$, $b_1=b_2=2$, and $l=d-1$, the resulting $\bm{R}^G(m_1,b_1,m_2,b_2,l)$ spans the entire set of $d$-dimensional orthogonal matrices. This hierarchical arrangement is referred to as the kaleidoscope hierarchy~\cite{dao2019kaleidoscope}. It is important to emphasize, however, that increased expressivity does not guarantee improved finetuning performance in practice; for example, full finetuning—despite its theoretical universality—often yields inferior outcomes. The crux for finetuning lies in balancing expressivity with regularization to promote generalization. By introducing structural priors, BOFT broadens the search space for finetuning parameters, thereby facilitating a more optimal compromise between expressivity and regularity.

\textbf{Spectral properties}. Compared to LoRA, orthogonal finetuning typically achieves superior spectral characteristics, owing to its exact preservation of the spectral norm of the original weight matrix $\bm{W}^0$. This invariance becomes apparent through singular value decomposition: for $\bm{W}^0=\bm{U}\bm{\Sigma}\bm{V}^\top$, where $\bm{U}$ and $\bm{V}$ are orthogonal matrices and $\bm{\Sigma}$ is a diagonal matrix containing the singular values, both OFT and BOFT operate by left-multiplying with an orthogonal matrix $\bm{R}$. The resulting adapted weights, $\bm{R}\bm{U}\bm{\Sigma}\bm{V}^\top$, retain the largest singular value—meaning the spectral norm of $\bm{W}^0$ remains unchanged. Such preservation of the spectral norm is known to contribute significantly to improved training stability and better generalization performance~\cite{miyato2018spectral,yoshida2017spectral}. We further discuss notable mathematical consequences in Appendix~\ref{app:sn_property}.

\textbf{Orthogonal finetuning as learning bilinear similarity}. The BOFT approach can be framed as optimizing a bilinear similarity function, specifically $\bm{w}^0_i\bm{R}\bm{x}$, where $\bm{w}^0_i$ denotes the $i$-th column (neuron) of the weight matrix $\bm{W}^0$. Under this lens, BOFT involves the estimation of the bilinear similarity matrix $\bm{R}$, with the additional constraint that $\bm{R}$ remains orthogonal. This ties BOFT conceptually to distance metric learning~\cite{xing2002distance} and the study of bilinear forms~\cite{roman2005advanced}.

\textbf{Inductive bias and generalization in BOFT}. The matrix $\bm{R}(m,b)$ utilized within BOFT is typically restricted to a structurally defined subset of the orthogonal group, effectively narrowing the hypothesis space and consequently inducing an inductive bias. We propose that this structured bias, as imparted by the butterfly factorization, can enhance generalization capabilities due to its shared organizational patterns with several foundational linear transforms~\cite{dao2019kaleidoscope}, such as the discrete Fourier, sine/cosine, and Hadamard transforms. In addition, the inherent sparsity present in the matrix factorization of BOFT is likely to provide supplementary implicit inductive biases~\cite{gunasekar2017implicit,li2018algorithmic,liu2019neural}.

\textbf{Contrast with butterfly-based sparse training}. Several studies~\cite{dao2019learning,dao2019kaleidoscope,chen2022pixelated,dao2022monarch} have explored sparse training by utilizing the butterfly parameterization. These works commonly concentrate on directly reparameterizing the weight matrices via the butterfly structure and training neural networks from the ground up. In particular, \cite{dao2022monarch} investigates finetuning by initially projecting pretrained weights onto a modified form of butterfly matrices, then adapting the projected parts for target tasks. In contrast, BOFT introduces an alternative finetuning paradigm, in which transformation of weights is accomplished using layer-shared weight matrices.

\input{rebuttal-results}

\section{Concluding Remarks and Limitations}

This paper presents Orthogonal Butterfly, a versatile and parameter-efficient approach for finetuning foundation models, grounded in the butterfly matrix framework. The central innovation is enhancing parameter-efficiency by representing a dense orthogonal matrix as the product of several sparse orthogonal matrices. To facilitate the identification of appropriate matrix decompositions, we introduce a graphical framework based on information transmission. Within this perspective, we observe that the butterfly structure offers a practical and effective solution for sparse orthogonal matrix factorization. We provide experimental evidence supporting the utility of BOFT in finetuning various large-scale models, including language models, vision models, and text-to-image generators. Our findings further support BOFT’s advantage as a broadly applicable method for model finetuning.

Nonetheless, BOFT has certain shortcomings. Due to the need to express the final orthogonal matrix as a product of multiple orthogonal matrices, BOFT incurs a modestly higher training time overhead compared to OFT. Strategies to decrease BOFT’s training runtime represent an avenue for future research. On the positive side, after finetuning, the orthogonal matrices learned by BOFT can be directly absorbed into the pretrained model, introducing no extra latency during inference. Additionally, it is still unclear whether the butterfly network constitutes the most optimal architecture for information transmission. Our framework draws a parallel to the field of computer networking, which offers extensive insights into the efficiency of various network topologies for information dissemination. We anticipate that adopting alternative, more efficient network structures could further advance the construction of dense orthogonal matrices.